\documentclass[11pt]{article}
\usepackage[T1]{fontenc}
\usepackage[utf8]{inputenc}
\usepackage{lmodern}
\usepackage[a4paper,margin=25mm]{geometry}
\usepackage{amsmath,amssymb,amsthm,mathtools}
\usepackage{booktabs,longtable,array}
\usepackage[unicode,hidelinks]{hyperref}
\usepackage{microtype}
\newcommand{\E}{\mathbb E}
\newcommand{\Pbb}{\mathbb P}
\newcommand{\TV}{\operatorname{TV}}
\newtheorem{theorem}{Theorem}[section]
\newtheorem{proposition}[theorem]{Proposition}
\newtheorem{lemma}[theorem]{Lemma}
\newtheorem{corollary}[theorem]{Corollary}
\theoremstyle{definition}
\newtheorem{counterexample}[theorem]{Counterexample}
\newtheorem{definition}[theorem]{Definition}
\theoremstyle{remark}
\newtheorem{remark}[theorem]{Remark}
\setlength{\emergencystretch}{3em}
\allowdisplaybreaks
\title{Round 3: Theorem and Proof Appendix\\
Decision-Relevant Uncertainty for Deterministic Boundary Learning}
\author{Research synthesis for Burak's MSc thesis}
\date{5 October 2026}
\begin{document}
\maketitle
\begin{abstract}
This standalone appendix proves the finite-pool, latent-observation,
model-perturbation, sequential, physics-transfer, discovery and geometric
statements used in the Round 3 report. It distinguishes a coherent
Bayesian experiment from an approximate refitting algorithm, and a
deterministic physical sign from a stochastic working likelihood.
PROVED means valid under the stated mathematical assumptions; it is not
a claim of research priority or an empirical result about SPH.
COUNTEREXAMPLE identifies an explicit failure of a proposed implication.
KNOWN/PRIOR ART identifies established principles or specializations.
\end{abstract}
\tableofcontents
\clearpage
\section*{Assumptions and relation to the previous rounds}
\addcontentsline{toc}{section}{Assumptions and relation to the previous rounds}
\paragraph{Fixed decision contract [KNOWN/PRIOR ART].}
Targets, candidate queries, loss and admissible actions are fixed when
comparing current and expected future risks. Binary variables are spins
in $\{-1,+1\}$ unless a construction explicitly uses bits in $\{0,1\}$.
Hamming weights are nonnegative and sum to one except where unit weights
are explicitly used; a common positive normalization does not change
competitive ratios. Conditioning is exact under the declared joint law
unless a result explicitly concerns an approximate updater. Expectations
over latent worlds and over observation noise must not be conflated.
\paragraph{Observation contract [PROVED definitions].}
A deterministic observation reveals the target sign. A logistic/probit
observation draws a Bernoulli label conditional on the latent field.
These define different experiments even when their current predictive
class probabilities happen to coincide. A stochastic working likelihood
does not establish physical randomness of a deterministic simulator.
\paragraph{Scope of inherited results [KNOWN/PRIOR ART within this project].}
The earlier $B/N$ sequential margin bound and the sharp $4/5$ result for
$N=3,B=2$ remain restricted to ideal updating, noiseless revelation,
equal-weight full-pool Hamming loss and the fixed original pool.
They are not claimed for geometric losses or the sensor examples below.
Their full proofs and rational certificates are in the accompanying
Round 1--2 reference appendix and original verification package.
The new common-sign-noise theorem below includes the one-step noiseless
$1/(N-1)$ theorem as its $\alpha=1$ case.
\paragraph{Geometry and transfer [PROVED definitions].}
Graph and normal-tube results state their coverage, derivative,
density and topology assumptions explicitly. Shrinkage results use an
exact normal-location experiment and a declared shift class, not an
unstated normal approximation to binary GP training.
\paragraph{Priority [KNOWN/PRIOR ART].}
Bayes risk reduction, SUR, Gaussian conditioning, simulation bounds,
adaptive-submodular greedy guarantees and normal-mean shrinkage have
substantial prior literature. The separate literature-priority table
identifies the closest sources and leaves priority of specific finite
constructions unresolved. No novelty follows from a bounded search.
% Appendix fragment: no preamble or document environment.
% The parent document defines theorem, proposition, lemma and counterexample.

\section{Binary observations, latent targets, and coherent decision value}

Throughout this appendix, a binary target or observation takes values in
$\{-1,+1\}$ unless a counterexample explicitly uses $0/1$ notation.
The target set and nonnegative weights $w_i$ are fixed before the query.
The loss is weighted Hamming loss
$L(t,a)=\sum_i w_i\mathbf{1}\{t_i\ne a_i\}$, and every coordinatewise
binary action is allowed. Write $W=\sum_i w_i$.

\begin{lemma}[PROVED: binary-observation identity]
\label{lem:binary-observation-r3}
For any coherent joint law of $(T,O)$, put
$m=\mathbb{E}T$ and $c=\mathbb{E}(TO)$. The minimum expected error after
observing $O$ is
\[
 r^O=\frac{1-\max\{|m|,|c|\}}{2}.
\]
Consequently, for target coordinates $T_i$ and candidate observation $O_j$,
\[
 V(j)=\frac12\sum_i w_i\bigl(|c_{ij}|-|m_i|\bigr)_+,
 \qquad m_i=\mathbb{E}T_i,\quad c_{ij}=\mathbb{E}(T_iO_j).
\]
\end{lemma}
\begin{proof}
There are exactly four deterministic maps from a binary observation to a
binary action: $+1,-1,O,-O$. Their correlations with $T$ are respectively
$m,-m,c,-c$, and their errors are one half of one minus that correlation.
Randomization cannot improve a linear expected loss. The best correlation
is therefore $\max\{|m|,|c|\}$. Before querying, the best correlation is
$|m|$. Subtract the two errors and sum coordinatewise, which is legitimate
because Hamming loss is additive and the action space is unrestricted.
\end{proof}

\begin{proposition}[KNOWN/PRIOR ART; PROVED: nonnegative coherent Bayes value]
\label{prop:nonnegative-r3}
Let a fixed decision problem have integrable loss $L(T,a)$ and action set
$\mathcal A$. If the same coherent law is used before and after an
observation $O$, the expected posterior Bayes risk is no larger than the
prior Bayes risk. This applies to any loss, provided the posterior action
minimizes that same loss over the same action set.
\end{proposition}
\begin{proof}
For every prequery action $a\in\mathcal A$, the constant posterior policy
that ignores $O$ is admissible. Hence
\[
 \mathbb{E}\!\left[\inf_{b\in\mathcal A}
          \mathbb{E}[L(T,b)\mid O]\right]
 \le \mathbb{E}\!\left[\mathbb{E}[L(T,a)\mid O]\right]
 =\mathbb{E}L(T,a).
\]
Take the infimum over $a$. For finite binary Hamming decisions, existence
and measurability of minimizers present no difficulty. Equivalently, for
$q_i=\mathbb{P}(T_i=1)$ and $q_i'=\mathbb{P}(T_i=1\mid O)$, the concavity
of $h(q)=\min(q,1-q)$ and the identity $\mathbb{E}q_i'=q_i$ imply
$\mathbb{E}h(q_i')\le h(q_i)$.
\end{proof}

\begin{proposition}[PROVED: Gaussian/probit target-observation moments]
\label{prop:gaussian-probit-r3}
Suppose $F\sim N(\mu,\Sigma)$ is the current working posterior,
$T_i=\operatorname{sign}F_i$, and
$O_j=\operatorname{sign}(F_j+\eta_j)$ with independent
$\eta_j\sim N(0,\tau^2)$, $\tau>0$. For $s_i^2=\Sigma_{ii}>0$,
$g_j^2=\Sigma_{jj}+\tau^2$, and
$\rho_{ij}=\Sigma_{ij}/(s_i g_j)$,
\begin{align*}
 m_i&=2\Phi(\mu_i/s_i)-1,\\
 c_{ij}
 &=4\Phi_2(\mu_i/s_i,\mu_j/g_j;\rho_{ij})
   -2\Phi(\mu_i/s_i)-2\Phi(\mu_j/g_j)+1.
\end{align*}
In particular, for a centered candidate with standard deviation $s$,
\[
 V_{\rm self}
 =\frac1\pi\arcsin\!\left(\frac{s}{\sqrt{s^2+\tau^2}}\right)
 =\frac1\pi\arctan(s/\tau).
\]
\end{proposition}
\begin{proof}
The pair $(F_i,F_j+\eta_j)$ is Gaussian with the displayed variances and
correlation. Its positive-positive orthant probability is the displayed
bivariate normal CDF, by changing both centered standardized signs.
For two binary signs, direct expansion gives
$\mathbb{E}(TO)=4\mathbb{P}(T=1,O=1)-2\mathbb{P}(T=1)
-2\mathbb{P}(O=1)+1$. This proves the general formulas.
At zero means, the Gaussian quadrant identity
$\Phi_2(0,0;\rho)=1/4+\arcsin(\rho)/(2\pi)$ yields
$m=0$ and $c=2\arcsin(\rho)/\pi$. Apply
Lemma~\ref{lem:binary-observation-r3}. The trigonometric equality follows
because $s,\tau>0$. These are exact statements for the current Gaussian
law and the stated probit experiment; they do not make a logistic
posterior or a fresh Laplace refit exact.
\end{proof}

\begin{counterexample}[COUNTEREXAMPLE; PROVED: logistic and probit margin ratios]
\label{ce:logistic-margin-r3}
For both a fixed logistic and a fixed probit observation channel, there
are two independent latent targets for which observed-label margin
uniquely selects a candidate whose coherent Hamming-value ratio tends
to zero.
\end{counterexample}
\begin{proof}
Let $S$ be a fair sign, let $U$ be an independent sign with
$\mathbb{E}U=2\delta>0$, and take
$F_a=\varepsilon S$, $F_b=M U$. Write the channel as
$\mathbb{E}[O\mid F=f]=a(f)$, with either
$a(f)=\tanh(f/2)$ or $a(f)=2\Phi(f/\tau)-1$.
Choose $M>0$ and $0<2\delta<a(M)$. Then
$\mathbb{E}O_a=0$, while $\mathbb{E}O_b=2\delta a(M)>0$,
so observed-label margin selects $a$ uniquely.
The target means are $m_a=0$, $m_b=2\delta$ and the self-moments are
$c_{aa}=a(\varepsilon)$ and $c_{bb}=a(M)$.
Independence gives
$c_{ba}=m_b\mathbb{E}O_a=0$ and
$c_{ab}=m_a\mathbb{E}O_b=0$, so cross-target gains vanish. Thus
\[
 \frac{V(a)}{V(b)}
 =\frac{a(\varepsilon)}{a(M)-2\delta}\longrightarrow0.
\]
For sufficiently small $\varepsilon$, $b$ is the optimal candidate.
\end{proof}

\begin{proposition}[PROVED: centered Gaussian logistic self-value]
If $F=sZ$ with $Z\sim N(0,1)$ and the observation channel is logistic,
then
\[
 V_{\rm self}(s)=\frac12\mathbb{E}\tanh(s|Z|/2)
 \sim \frac{s}{\sqrt{8\pi}}\qquad(s\downarrow0).
\]
The probit channel with $\tau^2=8/\pi$ has the same first-order limit.
\end{proposition}
\begin{proof}
The target mean is zero by symmetry. Its target-observation moment is
$\mathbb{E}[\operatorname{sign}(F)\tanh(F/2)]
=\mathbb{E}\tanh(s|Z|/2)$, proving the equality.
Since $\tanh x\le x$ for $x\ge0$, dominated convergence after division by
$s$ gives limit $\mathbb{E}|Z|/4=1/\sqrt{8\pi}$.
The probit formula follows from
$\arctan(s/\tau)/\pi\sim s/(\pi\tau)$.
\end{proof}

\begin{counterexample}[COUNTEREXAMPLE; PROVED: latent-scale gauge]
\label{ce:scale-gauge-r3}
Let $G$ be any random field with
$\mathbb{P}(G_i=0)=0$ at all finitely many target coordinates, and put
$F^{(c)}=cG$ for $c>0$. All sign laws and samplewise zero sets are
unchanged by $c$. Nevertheless, under any fixed bounded local channel
with response $a$ continuous at zero and $a(0)=0$, the one-query
Hamming values tend to zero as $c\downarrow0$.
\end{counterexample}
\begin{proof}
Positive scaling leaves every sign and zero set unchanged. Conditional
on $G$, the observation mean is $a(cG_j)$, so
\[
 \mathbb{E}[T_iO_j]
 =\mathbb{E}[\operatorname{sign}(G_i)a(cG_j)]
 \longrightarrow0
\]
by bounded dominated convergence. Lemma~\ref{lem:binary-observation-r3}
then gives $V_c(j)\to0$ for every candidate. In a Gaussian example,
means and covariance scale as $c\mu$ and $c^2\Sigma$; standardized means
$\mu_i/s_i$ remain fixed. In particular a centered sign has entropy one
bit for every $c>0$, even though its latent standard deviation vanishes.
When defined, the local boundary-displacement proxy
$s_i/|\nabla\mu_i|$ is also invariant. Under the different, noiseless
channel $O_j=T_j$, all query values are invariant rather than vanishing.
\end{proof}

\begin{theorem}[PROVED: sharp margin guarantee under common sign noise]
\label{thm:bsc-margin-r3}
Suppose the $N\ge2$ candidates are exactly the $N$ unit-weight targets.
For each candidate let $O_j=T_jB_j$, where $B_j$ is independent of the
entire target field and $\mathbb{E}B_j=\alpha\in(0,1]$ is common to all
candidates. Every observed-label margin maximizer $m$ satisfies
\[
 V(m)\ge\frac{1}{N-1}\max_k V(k).
\]
The constant is sharp for every fixed $\alpha>0$.
\end{theorem}
\begin{proof}
Put $a_i=|\mathbb{E}T_i|$ and $q_{ik}=|\mathbb{E}T_iT_k|$.
Since $\mathbb{E}O_j=\alpha\mathbb{E}T_j$, margin minimizes $a_j$.
The value identity becomes
$2V(k)=\sum_i(\alpha q_{ik}-a_i)_+$.
Let $f_b(x)=(b-x)_+$, $s_i=f_\alpha(a_i)$, and
$d_{ik}=f_{\alpha q_{ik}}(a_i)$. For $a_m\le a_k$ and $b\le\alpha$,
\[
 f_\alpha(a_m)-f_\alpha(a_k)
 \ge f_b(a_m)-f_b(a_k).
\]
Indeed the left and right sides are the lengths of the intersections of
$[a_m,a_k]$ with $[0,\alpha]$ and $[0,b]$, respectively.
Using $b=\alpha q_{mk}$ and symmetry of $q_{mk}$ gives
$s_k+d_{mk}\le s_m+d_{km}$.
Every remaining term satisfies $d_{ik}\le s_i\le s_m$, whence
\[
 2V(k)
 \le s_m+d_{km}+(N-2)s_m
 \le (N-1)\,2V(m).
\]
If $V(m)=0$, this inequality also proves all values are zero.

For sharpness when $N\ge3$, take one independent fair decoy sign and
$N-1$ identical signs of mean $2\eta>0$. Margin uniquely selects the
decoy. Its value is $\alpha/2$, whereas a cluster query has value
$(N-1)(\alpha-2\eta)/2$ when $2\eta<\alpha$. For sufficiently small
$\eta$, the cluster query is optimal, and the ratio approaches
$1/(N-1)$ as $\eta\downarrow0$. For $N=2$, the proved constant is one
and cannot be improved. Independence among the different noise variables
$B_j$ is not needed because only one observation is acquired.
\end{proof}

\begin{counterexample}[COUNTEREXAMPLE; PROVED: sign calibration misses amplitude]
\label{ce:amplitude-swap-r3}
For a fixed symmetric logistic or probit channel, identical full
latent-sign laws and identical observation marginals need not yield
identical optimal queries.
\end{counterexample}
\begin{proof}
Take independent centered Gaussian coordinates with standard deviations
$(1,\varepsilon)$ under $P$ and $(\varepsilon,1)$ under $Q$.
Under both laws the sign vector consists of two independent fair signs;
all sign marginals and sign-pair probabilities therefore agree exactly.
Symmetry also makes every observation marginal fair.
The self-value of the unit-scale coordinate is strictly positive,
whereas that of the $\varepsilon$-scale coordinate tends to zero by the
preceding propositions. Independence eliminates cross-target gains.
Consequently the optimal query swaps between $P$ and $Q$.
No score depending only on latent-sign joint probabilities, even when
supplemented with observation marginals, detects this discrepancy.
\end{proof}

\section{Transfer between working and comparison laws}

In this section $P$ and $Q$ concern the same binary target coordinates,
candidate observation actions, fixed weights, and loss. They may differ
in latent distribution or observation channel.

\begin{proposition}[PROVED: moment transfer for Bayes query values]
\label{prop:moment-transfer-r3}
Define
\[
 M=\sum_i w_i|m_i^P-m_i^Q|,
 \qquad C_j=\sum_i w_i|c_{ij}^P-c_{ij}^Q|.
\]
Then
$|V_P(j)-V_Q(j)|\le(M+C_j)/2$.
If $j_P$ and $j_Q$ maximize the corresponding values, then
\[
 V_P(j_P)-V_P(j_Q)
 \le M+\frac{C_{j_P}+C_{j_Q}}{2}
 \le M+\max_j C_j.
\]
\end{proposition}
\begin{proof}
The map $(m,c)\mapsto(|c|-|m|)_+$ is one-Lipschitz in each coordinate,
so apply the triangle inequality to the binary-observation formula.
For the second claim insert
$V_Q(j_P)-V_Q(j_Q)\le0$ between the two $P$-values and apply the
per-candidate bound to each remaining difference.
\end{proof}

\begin{theorem}[PROVED: tight terminal-policy transfer]
\label{thm:terminal-transfer-r3}
Let
\[
 D_j=\sum_i w_i
 \max\{|m_i^P-m_i^Q|,\ |c_{ij}^P-c_{ij}^Q|\}.
\]
Let $j_Q$ minimize terminal Bayes risk under $Q$, and after this query
use a $Q$-optimal decision rule $A^Q_{j_Q}(O_{j_Q})$.
Let $j_P$ minimize terminal Bayes risk under $P$, with terminal risk
$r_P^*(j_P)$. Then the complete query-and-decision policy satisfies
\[
 \mathbb{E}_P L(T,A^Q_{j_Q}(O_{j_Q}))-r_P^*(j_P)
 \le \frac{D_{j_Q}+D_{j_P}}{2}
 \le \max_j D_j.
\]
The universal leading constant one in the last bound cannot be reduced.
The same bound controls query-selection regret when the final decision
after $j_Q$ is instead optimized under $P$.
\end{theorem}
\begin{proof}
Fix a candidate $j$ and a deterministic coordinatewise rule.
Each coordinate rule is one of $+1,-1,O_j,-O_j$, so its correlation
with $T_i$ is one of $\pm m_i,\pm c_{ij}$.
Its error differs between $P$ and $Q$ by at most half the maximum
moment difference appearing in $D_j$. Summing gives
\[
 |\operatorname{risk}_P(A_j)-\operatorname{risk}_Q(A_j)|
 \le D_j/2.
\]
Taking minima also yields $|r_P^*(j)-r_Q^*(j)|\le D_j/2$.
Since $j_Q$ is $Q$-optimal,
\begin{align*}
 \operatorname{risk}_P(A^Q_{j_Q})
 &\le r_Q^*(j_Q)+D_{j_Q}/2\\
 &\le r_Q^*(j_P)+D_{j_Q}/2\\
 &\le r_P^*(j_P)+(D_{j_Q}+D_{j_P})/2.
\end{align*}
Optimizing the final action under $P$ can only decrease the left side.

To show sharpness, one target and one candidate suffice.
Choose $0<\epsilon<1/2$ and $0<\delta$ small.
Let $P$ have moments $(m,c,n)=(\epsilon,-\epsilon,0)$ and let $Q$
have $(m,c,n)=(0,\delta,0)$, where $n=\mathbb{E}O$.
Both are valid laws because
$\mathbb{P}(T=t,O=o)=(1+tm+on+toc)/4$ is nonnegative for these choices.
Under $Q$, the unique optimal rule is $A=O$.
Under $P$ its error is $(1+\epsilon)/2$, whereas optimal error is
$(1-\epsilon)/2$. Regret equals $\epsilon$, while
$D=\epsilon+\delta$. Their ratio tends to one as $\delta\downarrow0$.
\end{proof}

\begin{proposition}[PROVED: latent marginals, conditional structure, and channel]
\label{prop:channel-split-r3}
Assume local channels
$\mathbb{E}_R[O_j\mid F]=a_j^R(F_j)$ for $R\in\{P,Q\}$,
with $|a_j^R|\le1$.
Write $\mu_j^R=\operatorname{Law}_R(F_j)$ and
$h_{ij}^R(f)=\mathbb{E}_R[T_i\mid F_j=f]$.
Using total variation distance
$\operatorname{TV}(\mu,\nu)=\sup_A|\mu(A)-\nu(A)|$, one has
\begin{align*}
 |c_{ij}^P-c_{ij}^Q|
 \le{}&2\operatorname{TV}(\mu_j^P,\mu_j^Q)\\
 &+\mathbb{E}_Q|h_{ij}^P(F_j)-h_{ij}^Q(F_j)|\\
 &+\mathbb{E}_Q|a_j^P(F_j)-a_j^Q(F_j)|.
\end{align*}
The terms separate candidate latent marginal error, latent conditional
structure error, and observation channel error. Target marginal errors
remain explicit in the preceding transfer bounds.
\end{proposition}
\begin{proof}
The local channel assumption gives
$c_{ij}^R=\int h_{ij}^R a_j^R\,d\mu_j^R$.
Add and subtract $\int h_{ij}^P a_j^P\,d\mu_j^Q$ and then
$\int h_{ij}^Q a_j^P\,d\mu_j^Q$.
The first difference is bounded by twice total variation because the
integrand has absolute value at most one. The other two differences are
bounded by the displayed expectations, since $|h|,|a|\le1$.
One may choose any bounded version of $h^P$ on $Q$-only support;
the identity and bound remain valid, although the bound can then be
uninformative. The conditional-structure term is not asserted to be a
pure copula discrepancy.
\end{proof}

\begin{proposition}[PROVED: alternative marginal/dependence/channel split]
Let $\gamma_{ij}^R=\operatorname{Law}_R(F_i,F_j)$ and
$\kappa_{ij}^R=\gamma_{ij}^R-\mu_i^R\otimes\mu_j^R$.
Set $\delta_i=\operatorname{TV}(\mu_i^P,\mu_i^Q)$ and let
$\|\cdot\|_{\rm var}$ denote the full variation norm of a signed measure.
Under the preceding local-channel assumptions,
\begin{align*}
 |c_{ij}^P-c_{ij}^Q|
 \le{}&2\delta_i+2\delta_j+
       \|\kappa_{ij}^P-\kappa_{ij}^Q\|_{\rm var}\\
 &+\mathbb{E}_Q|a_j^P(F_j)-a_j^Q(F_j)|,
 \qquad |m_i^P-m_i^Q|\le2\delta_i.
\end{align*}
\end{proposition}
\begin{proof}
First keep the $P$ channel fixed. The integrand
$\operatorname{sign}(f_i)a_j^P(f_j)$ is bounded by one.
Decompose the pair-law difference into its product-marginal difference
and $\kappa_{ij}^P-\kappa_{ij}^Q$.
The product measures have total variation distance at most
$\delta_i+\delta_j$, as follows by adding and subtracting
$\mu_i^Q\otimes\mu_j^P$. Integration against their difference is
therefore bounded by $2\delta_i+2\delta_j$.
The dependence part is bounded by its full variation norm.
Finally change the channel under $Q$ and use its displayed expectation.
The marginal sign bound follows directly from the bounded sign function.
\end{proof}

\begin{proposition}[PROVED: algorithmic gain versus Bayes information]
\label{prop:updater-decomposition-r3}
For a fixed prequery action $a_0$ and algorithmic postquery rule
$A_j(O_j)$, define
\[
 b_P=\operatorname{risk}_P(a_0)-r_P^0,\qquad
 e_P(j)=\operatorname{risk}_P(A_j)-r_P^*(j).
\]
Both are nonnegative, and the algorithmic gain is exactly
\[
 G_P^A(j)=V_P(j)+b_P-e_P(j).
\]
Consequently, for any model-selected query $j_Q$,
\[
 \max_jG_P^A(j)-G_P^A(j_Q)
 =\max_j\{V_P(j)-e_P(j)\}
    -\{V_P(j_Q)-e_P(j_Q)\}.
\]
This residual does not, without additional controls, identify
information absent from the working model.
\end{proposition}
\begin{proof}
By definition,
$\operatorname{risk}_P(a_0)=r_P^0+b_P$ and
$\operatorname{risk}_P(A_j)=r_P^*(j)+e_P(j)$.
Subtract and use $V_P(j)=r_P^0-r_P^*(j)$.
The second identity follows because $b_P$ is independent of $j$.
Even if the working law equals $P$, different updater excess risks
$e_P(j)$ can leave positive algorithmic regret.
If $P$ is conditioned on a fully known realized target field, then
$r_P^0=r_P^*(j)=V_P(j)=0$; nevertheless a restricted learner may have
positive gains because its own errors change. Such fixed-truth gains
are algorithmic utilities, not positive Bayes information under the
omniscient conditional law.
\end{proof}

\section{Coherent negative values for other risk functionals}

The next examples use $T_1,T_2\in\{0,1\}$.
A plug-in edge prediction is the exclusive-or of the two marginal
mode predictions. Its edge risk is the probability that this predicted
cut differs from the true cut. Its Dice ratio is
\[
 R_D=\frac{\mathbb{E}W_{\rm err}}
           {\mathbb{E}W_{\rm true}+W_{\rm pred}}.
\]
Neither operation is automatically the minimum posterior expectation
of one fixed loss, so Proposition~\ref{prop:nonnegative-r3} does not
automatically apply.

\begin{counterexample}[COUNTEREXAMPLE; PROVED: negative coherent plug-in edge gain]
Assign masses $(2,2,2,3)/9$ to states $(00,01,10,11)$ and observe
$O=\mathbf{1}\{(T_1,T_2)=(0,1)\}$. Exact Bayesian conditioning gives
an expected plug-in edge-risk reduction of $-1/9$.
\end{counterexample}
\begin{proof}
Before observing, both marginal modes are one, so the edge is predicted
uncut. Its error is the true-cut probability $4/9$.
The branch $O=1$ identifies state $01$ exactly and has zero error.
On $O=0$, whose probability is $7/9$, the conditional masses on
$(00,10,11)$ are $(2,2,3)/7$. The marginal probabilities of one are
$5/7$ and $3/7$, so the plug-in modes are $(1,0)$ and predict a cut.
Its conditional error is the no-cut probability $5/7$.
Expected future error is $(7/9)(5/7)=5/9$, so the reduction is
$4/9-5/9=-1/9$. All marginal decisions are strict, and every update
is exact conditioning.
\end{proof}

\begin{counterexample}[COUNTEREXAMPLE; PROVED: negative coherent Dice-ratio gain]
Assign masses $(1,1,2,1)/5$ to $(00,01,10,11)$ and observe the same
$O=\mathbf{1}\{(T_1,T_2)=(0,1)\}$. Exact conditioning gives an expected
Dice-ratio reduction of $-1/60$.
\end{counterexample}
\begin{proof}
The prior marginal modes are $(1,0)$, the true-cut probability is $3/5$,
and the ratio is
$R_D=(2/5)/(1+3/5)=1/4$.
The $O=1$ branch has zero error and zero ratio.
On $O=0$, of probability $4/5$, conditional masses on $(00,10,11)$
are $(1,2,1)/4$. Marginal modes remain strictly $(1,0)$.
The true-cut probability is $1/2$, so
$R_D=(1/2)/(1+1/2)=1/3$.
Expected future ratio is $(4/5)(1/3)=4/15$ and the reduction is
$1/4-4/15=-1/60$.
Thus a negative ratio reduction by itself does not demonstrate an
incoherent posterior. A separately measured failure of the posterior
martingale identity is a different diagnostic.
\end{proof}


% Round 3 proof fragment. No preamble or document environment.
% Requires amsmath/amssymb/amsthm and theorem/proposition/lemma/counterexample
% environments; the parent document supplies \E and \TV. Links use hyperref.
% No thesis experiments were run and the repository was read-only.

\section{Sequential planning under a misspecified law}

\paragraph{Scope [KNOWN/PRIOR ART].}
Throughout this section, the target, feasible queries, horizon, and loss are
common to the truth law $P$ and planning law $Q$.  The existing Round~2
$B/N$ guarantee and the sharp $N=3,B=2$ ratio $4/5$ concern noiseless
full-pool Hamming loss and are retained without rederivation.  The sensor
counterexample below has a disjoint target and a noisy observation and is
outside that scope.

Let $H_t=(A_0,O_0,\ldots,A_{t-1},O_{t-1})$ be the history before query $t$.
A policy $\pi$ maps histories to distributions over feasible actions.
Its internal randomness is independent of the underlying state.  Its
terminal decision is $d(H_B)$, and $0\leq L(d,T)\leq1$ is the common loss
for a fixed target $T$.  Write $R_M(\pi,d)=\E_M^\pi L(d(H_B),T)$ for
$M\in\{P,Q\}$.  Evaluating a fixed policy under the two laws means using
the same action and terminal-decision maps, even when these maps were
computed using $Q$.

\begin{theorem}[PROVED: a history-dependent simulation bound]
Define
\begin{align*}
 K_{M,t}(\cdot\mid h,a)&=M(O_t\in\cdot\mid H_t=h,A_t=a),\\
 \epsilon_t(h,a)&=\TV\bigl(K_{P,t}(\cdot\mid h,a),
                              K_{Q,t}(\cdot\mid h,a)\bigr),\\
 \delta(h)&=\TV\bigl(P(T\in\cdot\mid H_B=h),
                         Q(T\in\cdot\mid H_B=h)\bigr),\\
 E_\pi&=\sum_{t=0}^{B-1}\E_P^\pi\epsilon_t(H_t,A_t),
 &D_\pi&=\E_P^\pi\delta(H_B),\\
 \Gamma_\pi&=\min\{1,E_\pi+D_\pi\}.
\end{align*}
Use $\TV(P,Q)=\sup_A|P(A)-Q(A)|$, one half of the $L^1$ distance.
Specify versions of the $Q$ kernels and posterior at every $P$-reachable
history, including $Q$-null histories.  Then every terminal rule satisfies
\begin{equation}
 |R_P(\pi,d)-R_Q(\pi,d)|\leq\Gamma_\pi.
 \label{eq:r3-sequential-simulation}
\end{equation}
\end{theorem}

\begin{proof}
Let $\mu_t$ and $\nu_t$ be the $P$ and $Q$ laws of $H_t$ under the fixed
policy.  Their initial histories agree.  Incorporate the common action
kernel into $K_{P,t}$ and $K_{Q,t}$, and apply contraction of total variation
under a Markov kernel:
\begin{align*}
 \TV(\mu_{t+1},\nu_{t+1})
 &\leq \TV(\mu_tK_{P,t},\mu_tK_{Q,t})
       +\TV(\mu_tK_{Q,t},\nu_tK_{Q,t})\\
 &\leq\E_{\mu_t,\pi}\epsilon_t(H_t,A_t)
       +\TV(\mu_t,\nu_t).
\end{align*}
Induction gives $\TV(\mu_B,\nu_B)\leq E_\pi$.
Put $r_M(h)=\E_M[L(d(h),T)\mid H_B=h]$.  The terminal posterior comparison
gives $|r_P(h)-r_Q(h)|\leq\delta(h)$, and $0\leq r_Q\leq1$.  Therefore
\begin{align*}
 |\E_{\mu_B}r_P-\E_{\nu_B}r_Q|
 &\leq\E_{\mu_B}|r_P-r_Q|
          +|\E_{\mu_B}r_Q-\E_{\nu_B}r_Q|\\
 &\leq D_\pi+\TV(\mu_B,\nu_B)\leq D_\pi+E_\pi.
\end{align*}
The difference of two risks in $[0,1]$ is at most one.  Randomized terminal
decisions can be averaged in $r_M$, so the same proof applies.
\end{proof}

\begin{proposition}[PROVED: only terminal marginals are needed for Hamming]
For binary targets $T_i$, fixed weights $w_i\geq0$, $\sum_iw_i=1$, and
loss $L(d,T)=\sum_iw_i\mathbf{1}\{d_i\ne T_i\}$, the preceding theorem
remains true on replacing $\delta$ by
\begin{equation}
 \delta_H(h)=\sum_iw_i
 \left|P(T_i=1\mid h)-Q(T_i=1\mid h)\right|.
 \label{eq:r3-terminal-marginals}
\end{equation}
\end{proposition}
\begin{proof}
For a fixed binary decision the error probability is either $p_i$ or
$1-p_i$, whose change is bounded by $|p_{P,i}-p_{Q,i}|$.  Sum over $i$ and
repeat the proof above.  The required marginals are those conditional on
the terminal history, not merely those before acquisition.
\end{proof}

\begin{theorem}[PROVED: regret of a planning-law-optimal policy]
Suppose $(\pi_Q,d_Q)$ is $\beta$-optimal under $Q$ in a common class of
policies and terminal decisions, and $(\pi_P,d_P)$ is optimal under $P$.
Then
\begin{equation}
 R_P(\pi_Q,d_Q)-R_P(\pi_P,d_P)
 \leq\beta+\Gamma_{\pi_Q}+\Gamma_{\pi_P}.
 \label{eq:r3-sequential-regret}
\end{equation}
The right side can be capped by one.  The weighted-Hamming version uses
\eqref{eq:r3-terminal-marginals}.
\end{theorem}
\begin{proof}
Add and subtract $R_Q(\pi_Q,d_Q)$ and $R_Q(\pi_P,d_P)$.  The two
cross-law differences are bounded by
\eqref{eq:r3-sequential-simulation}; $\beta$-optimality bounds the remaining
within-$Q$ difference.  This also proves the result with infima and
arbitrarily small optimization errors if an optimizer is not attained.
\end{proof}

\paragraph{Interpretation [PROVED].}
The discrepancies are weighted by histories reached under $P$ by each
compared policy.  Checking only histories generated by an old margin
policy does not bound a new policy's error when it visits different
histories.  The $P$-optimal-policy term is not identifiable from a single
realized campaign without further information.

\begin{proposition}[PROVED: uniform coupling and relative-entropy versions]
If $\epsilon_t(h,a)\leq e_t\in[0,1]$ and full target-posterior
$\delta(h)\leq d\in[0,1]$ uniformly, then
\begin{equation}
 |R_P(\pi,d_0)-R_Q(\pi,d_0)|
 \leq1-(1-d)\prod_{t=0}^{B-1}(1-e_t)
 \leq d+\sum_{t=0}^{B-1}e_t
 \label{eq:r3-coupling}
\end{equation}
for any terminal rule $d_0$.
Under absolute continuity, define
\begin{align*}
 J_\pi={}&\sum_{t=0}^{B-1}\E_P^\pi
 \operatorname{KL}\bigl(K_{P,t}(\cdot\mid H_t,A_t)
                    \Vert K_{Q,t}(\cdot\mid H_t,A_t)\bigr)\\
 &+\E_P^\pi\operatorname{KL}\bigl(P(T\mid H_B)\Vert Q(T\mid H_B)\bigr).
\end{align*}
Then $|R_P(\pi,d_0)-R_Q(\pi,d_0)|\leq\sqrt{J_\pi/2}$, and the regret
is at most $\beta+\sqrt{J_{\pi_Q}/2}+\sqrt{J_{\pi_P}/2}$.
\end{proposition}
\begin{proof}
While histories coincide, couple actions with identical policy randomness
and observations by a maximal coupling.  The probability of remaining
matched through the observations is at least $\prod_t(1-e_t)$.  Conditional
on a matched terminal history, a maximal target coupling succeeds with
probability at least $1-d$.  The probability of any mismatch bounds the
difference of expectations of a common $[0,1]$ loss.  For the second claim,
the chain rule gives
$\operatorname{KL}(P^\pi_{H_B,T}\Vert Q^\pi_{H_B,T})=J_\pi$ because the
common policy action factors cancel.  Pinsker's inequality and the
regret decomposition complete the proof.  Without absolute continuity
the corresponding infinite bound is uninformative, not false.
\end{proof}

\paragraph{Prior art [KNOWN/PRIOR ART].}
The simulation-lemma strategy is classical; see Kearns and Singh,
\emph{Near-Optimal Reinforcement Learning in Polynomial Time}, Machine
Learning 49 (2002), 209--232,
\href{https://people.eecs.berkeley.edu/~pabbeel/cs287-fa09/readings/KearnsSingh_E3_2002.pdf}{primary paper}.
The display above separates the observation kernel from the terminal
latent-target posterior.  Its loss is terminal and normalized; cumulative
losses require summing prefix bounds and can introduce another horizon
factor.

\paragraph{Experiment row [CONJECTURE: test proposal].}
Prediction: small errors in both objects prevent large regret.
Model/data: a frozen finite sensor law or small discretized Gaussian world,
with exact history enumeration and common feasible actions.
Statistic: exact regret and the signed slack in
\eqref{eq:r3-sequential-regret}, together with each occupancy-weighted term.
Falsifier: negative certified slack.  Real observed labels alone do not
reveal the latent-target posterior or all counterfactual truth kernels.

\subsection{One-step correctness and long-horizon failure}

\begin{counterexample}[COUNTEREXAMPLE: parity defeats exact myopic value]
Let $A,B$ be independent fair signs, $T=AB$, and $E$ independent with
$P(E=1)=1/2+\eta$, $0<\eta<1/2$.  Queries reveal $A$, $B$, or $C=TE$,
each at most once.  The sole target is $T$, the loss is classification
error, and the budget is two.  Exact true-law myopic value has total gain
$\eta$, while the optimal two-query gain is $1/2$.
\end{counterexample}
\begin{proof}
$A$ is independent of $(T,E)$ because $B$ is fair; similarly $B$ is
independent of $(T,E)$.  Therefore the initial values of $A,B$ are zero.
Query $C$ predicts $T$ with accuracy $1/2+\eta$, so its value is $\eta$
and myopic value selects it strictly.  Either remaining query is
independent of $(T,C)$, hence adds no value.  Querying $A,B$ reveals $T$
exactly, giving gain $1/2$.  The gain ratio is $2\eta\to0$.
\end{proof}

\begin{counterexample}[COUNTEREXAMPLE: observation kernels and initial pairs can agree]
Keep the preceding truth law $P$.  Let $Q$ instead make $T,A,B$ independent
fair signs while preserving $C=TE$ and the same noise law.  All initial
pairs among $T,A,B,C$ agree under $P,Q$, and every adaptive observation
transcript of length at most two has the same law.  Nevertheless a
$Q$-optimal policy has $P$ risk $1/2-\eta$, while the $P$ optimum has
risk zero.
\end{counterexample}
\begin{proof}
Under each law, every distinct sensor pair among $A,B,C$ is independent
fair.  Thus the first observation is fair and, given its realized value,
either unqueried second observation is fair; this establishes transcript
equality even with adaptive action selection.  The pair $(T,C)$ has
agreement probability $1/2+\eta$ in both laws; $(T,A)$ and $(T,B)$ are
independent fair in both.  Under $Q$, $A,B$ together contain no target
information, so an optimal policy includes $C$.  Under $P$, their product
is the target.  After $A,B$, the $P$ target posterior is a point mass at
$AB$, while the $Q$ posterior is fair, at total variation $1/2$.  All
observation-kernel errors through budget two are zero, but this terminal
posterior error is nonzero.  The stated risks follow.
\end{proof}

\paragraph{Scope [PROVED].}
The example proves that initial target-observation pairwise correctness
does not determine multi-observation target posteriors.  It also proves
that the terminal-posterior term cannot be dropped from the simulation
bound.  It does not contradict noiseless full-pool Hamming guarantees.

\paragraph{Experiment row [CONJECTURE: test proposal].}
Use the finite sign worlds with rational $\eta=1/10$.  Enumerate all
one-step values and two-query trees, all initial pair laws, and the final
risks.  The exact predictions are gains $1/10$ and $1/2$, zero initial
pair discrepancy, and terminal discrepancy $1/2$ after $A,B$.  Any
different value falsifies the implementation of the construction.

\subsection{A sufficient condition for a greedy horizon guarantee}

\begin{theorem}[KNOWN/PRIOR ART: adaptive submodularity]
Let a nonnegative realization utility $f(S,\phi)$ be normalized by
$f(\varnothing,\phi)=0$.  Assume adaptive monotonicity and adaptive
submodularity under the actual common law: conditional expected gains
$\Delta(e\mid\psi)$ are nonnegative and
$\Delta(e\mid\psi)\geq\Delta(e\mid\psi')$ whenever $\psi\subseteq\psi'$
and $e$ is still unqueried.  For unit costs and budget $B$, exact adaptive
greedy has expected gain at least
\[
 \left[1-\left(1-\frac1B\right)^B\right]\operatorname{OPT}
 \geq(1-e^{-1})\operatorname{OPT}.
\]
\end{theorem}
\begin{proof}
Let $G_t$ denote expected greedy gain after $t$ selections.  At its current
partial history, adaptive submodularity bounds the expected incremental
value of any $B$-query continuation by $B$ times the best current marginal
gain.  Adaptive monotonicity allows comparison with concatenating the
current policy and the optimal $B$-query policy, omitting already known
queries and ignoring the first policy's extra information when simulating
the second.  Consequently
$\operatorname{OPT}-G_t\leq B(G_{t+1}-G_t)$.
Rearrange and iterate from $G_0=0$ to obtain the displayed bound.
\end{proof}

\paragraph{Attribution and limits [KNOWN/PRIOR ART; PROVED].}
See Golovin and Krause, JAIR 42 (2011),
\href{https://arxiv.org/pdf/1003.3967}{Theorem 5 in the primary PDF}.
These are law-specific, objective-specific premises, not properties of
all Bayesian acquisition.  The parity example violates adaptive
submodularity: $\Delta(B\mid\varnothing)=0$ but
$\Delta(B\mid A=a)=1/2$.  A result for version-space elimination does not
automatically give a result for Hamming, NSD, DC-BD, or q20.

\begin{proposition}[PROVED: an adaptive-modular Hamming special case]
For independent binary targets with noiseless self-observation and fixed
Hamming weights, a query of target $i$ has fixed gain
$w_i\min(p_i,1-p_i)$, independent of previous observations.  Selecting the
largest $B$ such values is optimal; for equal weights this is ideal margin.
\end{proposition}
\begin{proof}
Observation removes target $i$'s own error and independence leaves every
other target posterior unchanged.  Gains therefore add without history
dependence, and maximizing their sum means taking the largest values.
\end{proof}

\paragraph{Experiment row [CONJECTURE: test proposal].}
Enumerate all nested reachable histories in a finite independent-label
positive control and the parity negative control.  Record
$\max_{\psi\subseteq\psi',e}[\Delta(e\mid\psi')-\Delta(e\mid\psi)]$
and the greedy-to-optimal ratio.  A positive value falsifies the
submodularity premise; a below-bound ratio with a fully verified premise
falsifies the computation.  A small sample of histories can find a
violation but cannot certify the universal premise.

\section{Physics as transferable order and uncertain level}

\subsection{An exact normal experiment for campaign-threshold borrowing}

Let the monotone score threshold in a new campaign be
$\theta=\theta_0+\Delta$, where $\theta_0$ is an OLD or external anchor.
At a fixed paid budget suppose current-campaign information is exactly
\[
 Z=\theta+\varepsilon,\qquad \varepsilon\sim N(0,v),\quad v>0,
\]
with noise independent of the campaign shift.  This is an exact
normal-location model.  Approximating a binary-label threshold estimate
by such a $Z$ is a separate assumption.  Consider
$\widehat\theta_w=\theta_0+w(Z-\theta_0)$, $0\leq w\leq1$.
Here $w=0$ fixes the OLD level; $w=1$ discards the OLD level while retaining
the score's direction; intermediate values weaken level borrowing.

\begin{proposition}[PROVED: exact bias--variance and crossover formulas]
For a fixed shift $\Delta$,
\begin{equation}
 R_\Delta(w)=\E(\widehat\theta_w-\theta)^2
           =(1-w)^2\Delta^2+w^2v.
 \label{eq:r3-normal-risk}
\end{equation}
The anchor has higher risk than current-only estimation exactly when
$\Delta^2>v$.  For $w<1$, borrowing has higher risk than current-only
exactly when $\Delta^2>v(1+w)/(1-w)$.  If $w<w'\leq1$, increasing the
update weight improves risk exactly when
\[
 \Delta^2>v\frac{w+w'}{2-w-w'}.
\]
\end{proposition}
\begin{proof}
$\widehat\theta_w-\theta=-(1-w)\Delta+w\varepsilon$; the cross term has
zero expectation, proving \eqref{eq:r3-normal-risk}.  Subtract $v$ or
$R_\Delta(w)$ from the respective comparison risks and factor:
\begin{align*}
 R_\Delta(w)-v
 &=(1-w)\{(1-w)\Delta^2-(1+w)v\},\\
 R_\Delta(w')-R_\Delta(w)
 &=(w'-w)\{(w+w')v-(2-w-w')\Delta^2\}.
\end{align*}
The factors' signs give the claimed thresholds.
\end{proof}

\begin{theorem}[PROVED: bounded-shift minimax among affine rules]
For $|\Delta|\leq D$,
\begin{align}
 w_D&=\frac{D^2}{D^2+v},\label{eq:r3-shrink-weight}\\
 \inf_{0\leq w\leq1}\sup_{|\Delta|\leq D}R_\Delta(w)
 &=\frac{D^2v}{D^2+v}.\label{eq:r3-affine-minimax}
\end{align}
For $D>0$ this is strictly below both endpoint worst-case risks $D^2,v$.
It is not an unrestricted bounded-parameter minimax claim.
\end{theorem}
\begin{proof}
The supremum in \eqref{eq:r3-normal-risk} is
$(1-w)^2D^2+w^2v$.  Its unique minimum follows by differentiation, yielding
\eqref{eq:r3-shrink-weight}--\eqref{eq:r3-affine-minimax}.
For the final distinction, project the affine estimator onto the interval
$[\theta_0-D,\theta_0+D]$.  Projection weakly decreases squared distance to
every parameter in that interval and strictly decreases expected loss
on positive-probability Gaussian tails when $D,v>0$.  Continuity on the
compact parameter interval also makes the worst-case risk strictly
smaller than that of the unclipped affine rule.  Therefore the latter
cannot be unrestricted minimax on the bounded parameter interval.
\end{proof}

\begin{theorem}[KNOWN/PRIOR ART: moment-class robust Bayes, specialized to campaigns]
Let
\[
 \mathcal G_D=\{\nu:\E_\nu\Delta^2\leq D^2\}
\]
be an ambiguity set of shift distributions with independent
$N(0,v)$ observation noise.  Among all measurable estimators,
\begin{equation}
 \inf_a\sup_{\nu\in\mathcal G_D}
 \E_{\nu,\varepsilon}\{a(Z)-\theta\}^2
 =\frac{D^2v}{D^2+v}.
 \label{eq:r3-moment-minimax}
\end{equation}
The affine rule with weight $w_D$ attains the value, and
$\nu_*=N(0,D^2)$ is least favorable.
\end{theorem}
\begin{proof}
Integrating \eqref{eq:r3-normal-risk} over any $\nu\in\mathcal G_D$ shows
that the affine estimator's risk is at most
$(1-w)^2D^2+w^2v$.  At $w_D$ this proves the upper bound.
For a matching lower bound, $\nu_*=N(0,D^2)$ belongs to the ambiguity set.
Completing the square gives
\[
 \Delta\mid Z=z\sim N\left(
 \frac{D^2}{D^2+v}(z-\theta_0),\frac{D^2v}{D^2+v}\right).
\]
The posterior mean minimizes conditional squared loss among all
measurable estimators.  Its Bayes risk is the constant posterior variance
$D^2v/(D^2+v)$.  Every estimator's worst-prior risk is at least its risk
under this admissible prior, hence at least that Bayes risk.  The affine
upper bound attains it, proving the saddle-point statement.  If $D=0$,
the ambiguity set concentrates at zero and the conclusion is immediate.
\end{proof}

\paragraph{Attribution and distinction [KNOWN/PRIOR ART].}
This is the classical univariate second-moment bound with a Gaussian
least favorable prior, not a new minimax result; see Johnstone,
\emph{Gaussian Estimation: Sequence and Wavelet Models}, subsection on
univariate moment bounds,
\href{https://imjohnstone.su.domains/GE_08_09_17.pdf}{author's primary PDF}.
The campaign interpretation is the specialization used here.
A Gaussian prior with variance $D^2$ has mass outside $[-D,D]$, so this
proof cannot establish unrestricted minimaxity for the bounded-shift
problem.  The two problems happen to share the same affine weight.

\begin{proposition}[PROVED: no bounded shift assumption, no finite worst-case borrowing]
Over all fixed $\Delta\in\mathbb R$, the current-only estimator $Z$ is
minimax among all measurable estimators with risk $v$.  Every fixed
affine rule with $w<1$ has infinite worst-case risk.
\end{proposition}
\begin{proof}
$Z$ has constant risk $v$.  For a lower bound, use normal priors of
variance $M^2$.  Their optimal Bayes risks are $M^2v/(M^2+v)\to v$;
worst-parameter risk bounds every such Bayes risk from above.
Equation~\eqref{eq:r3-normal-risk} diverges as $|\Delta|\to\infty$ when
$w<1$.
\end{proof}

\paragraph{Design limitation [PROVED].}
If a design changes only $v$, the robust risk $D^2v/(D^2+v)$ is increasing
in $v$.  Thus the design ordering still minimizes variance.  The theorem
proves a borrowing/inference result, not improved acquisition.  A design
claim needs a model in which bias or information geometry also changes
with the design.

\paragraph{Experiment row [CONJECTURE: test proposal].}
Freeze $v,D$ and a shift grid before generating normal-location data.
Compare squared-error means with \eqref{eq:r3-normal-risk}; report the
maximum over the prespecified shifts and the clipped-affine control.
The prediction is the exact quadratic and the displayed crossover.
Systematic disagreement in this exact experiment falsifies the
implementation.  Failure of a separately tested normal approximation in
binary-label data rejects applicability rather than the normal theorem.

\subsection{An exact order-only threshold model}

\begin{theorem}[PROVED; KNOWN binary-search mechanism: finite-pool order transfer]
Let $s_1<\cdots<s_N$ and $y_i=\mathbf{1}\{i>k\}$ for
$k\in\{0,\ldots,N\}$.  A fixed OLD threshold $k_0$ has normalized
full-pool Hamming error $|k-k_0|/N$.  An order-only learner that repeatedly
bisects the feasible threshold ranks leaves at most
\[
 m_B=\left\lceil\frac{N+1}{2^B}\right\rceil
\]
possible ranks after $B$ noiseless labels.  Predicting a median feasible
rank has worst-case error at most $\lfloor m_B/2\rfloor/N$.  Exact
identification is possible and requires
$\lceil\log_2(N+1)\rceil$ queries in the worst case.
\end{theorem}
\begin{proof}
Querying $i$ reveals whether $k<i$, which cuts the ordered feasible set
into two contiguous parts.  A balanced cut reduces its cardinality by
half up to a ceiling, giving $m_B$.  A set of $m$ consecutive integer
ranks spans $m-1$ positions; a median is at most
$\lceil(m-1)/2\rceil=\lfloor m/2\rfloor$ positions from either endpoint.
Two threshold vectors at ranks $k,k'$ differ on exactly $|k-k'|$ labels.
For exact identification a binary decision tree of depth $B$ has at most
$2^B$ leaves and must distinguish $N+1$ possible vectors.  Balanced search
attains that lower bound.
\end{proof}

\begin{counterexample}[COUNTEREXAMPLE: in-domain efficiency versus transfer worst case]
Take $N=7$, $k_0=2$, and $B=1$.  The order-only rule queries $i=4$ and
predicts the lower median of the remaining feasible threshold ranks.
At the unchanged campaign $k=2$, its estimate is $1$ and its error $1/7$,
whereas the correct anchor has zero error.  Across all transferred ranks
$k\in\{0,\ldots,7\}$, its worst-case error is $2/7$, while the anchor's
worst-case error is $5/7$.
\end{counterexample}
\begin{proof}
When $k<4$, the remaining ranks are $0,1,2,3$ and the estimate is $1$;
when $k\geq4$, they are $4,5,6,7$ and the estimate is $5$.  The maximum
distance to those estimates is two.  The maximum distance from the fixed
anchor $2$ is five, attained at $k=7$.
\end{proof}

\paragraph{Prior art and scope [KNOWN/PRIOR ART; PROVED].}
See Dasgupta, \emph{Analysis of a greedy active learning strategy},
NIPS 2004 / author version 2005,
\href{https://cseweb.ucsd.edu/~dasgupta/papers/greedy.pdf}{introduction and Figure 1}.
Strictly increasing transformations of the score preserve the entire
search problem.  The guarantee protects against changed level while
requiring an unchanged noiseless monotone order.  It does not survive
arbitrary inversions, nonmonotone boundaries, or noise without additional
analysis.  The observed overlap in real log-$h$ label ranges prevents
claiming that exact realizability has already been established for SPH.

\paragraph{Experiment row [CONJECTURE: test proposal].}
Enumerate all threshold ranks on a fixed sorted synthetic pool.  Record
feasible-set cardinality and maximum Hamming error by budget; recover
the explicit $N=7$ contrast.  Any bound violation in the exact noiseless
model falsifies the calculation.  On real labels, the inversion count
$\sum_{i<j}\mathbf{1}\{y_i=1,y_j=0\}$ tests the noiseless single-threshold
premise; a positive count rejects that premise, not every noisy monotone
association.

\subsection{Regularization is not automatically robust Bayes}

\begin{proposition}[PROVED: attenuating an affine mean preserves its threshold]
For $m(s)=b_0+b_1s$, $b_1\ne0$, and $\lambda>0$, multiplying both
coefficients by $\lambda$ leaves the zero crossing $-b_0/b_1$ unchanged.
A free campaign intercept $c$ changes the crossing to
$-(b_0+c)/b_1$.
\end{proposition}
\begin{proof}
Solve $\lambda(b_0+b_1s)=0$ and $b_0+b_1s+c=0$, respectively.
\end{proof}

\paragraph{Interpretation [PROVED; KNOWN/PRIOR ART].}
Mean attenuation can change confidence and interaction with a fitted
residual, but it cannot by itself repair a pure threshold shift.
Penalized logistic coefficients admit a MAP interpretation only after
the loss scaling, standardization, intercept convention, and prior are
specified.  A choice such as $C=1$ alone supplies no shift radius,
ambiguity set, least favorable prior, or minimax guarantee.

\paragraph{Repository mapping [NUMERICALLY CHECKED: source audit].}
The original M3 estimates a physics trend from its supplied training
labels and then fixes that trend in the residual-GP stage; strict
OLD-to-NEW transfer fixes an OLD-estimated level.  Week~16 instead uses
a standardized-score, $C=1$ logistic trend fitted on the currently
revealed labels, with zero-mean fallback for one class.  These differ
in campaign adaptation and regularization.  A plug-in two-stage fit also
omits posterior uncertainty in the estimated trend.  Reusing labels in
both stages may define an empirical-Bayes/composite procedure, but is
not equivalent to joint hierarchical Bayes with a single label
likelihood and uncertainty propagated through both components.

\paragraph{Experiment row [CONJECTURE: proposed mechanism test].}
Use a frozen logistic threshold-shift family with matched paid labels
and acquisition paths.  Separate intercept adaptation, slope penalty,
and residual flexibility.  Report threshold bias, variance, final risk,
and uncertainty.  The conjecture is that appropriately weakened level
borrowing improves shifted-campaign risk while preserving useful order.
A claim that positive whole-mean attenuation moves the zero crossing is
algebraically false; a benefit disappearing after matched controls does
not support a regularization-specific mechanism.

\section{What OLD-only calibration can identify}

\begin{theorem}[PROVED: an OLD-only threshold-transfer impossibility]
Suppose the OLD labeled-data law and the unlabeled target law are
identical in two worlds.  Let target $X\sim\operatorname{Uniform}[0,1]$
in both worlds, with labels $\mathbf{1}\{X>\theta_0\}$ and
$\mathbf{1}\{X>\theta_1\}$, respectively, where
$0\leq\theta_0<\theta_1\leq1$.  A learner using OLD labels and any
quantity of unlabeled target data, but no target labels, has risks
\[
 R_0+R_1\geq\theta_1-\theta_0,
 \qquad \max(R_0,R_1)\geq\frac{\theta_1-\theta_0}{2}.
\]
\end{theorem}
\begin{proof}
The learner's input and hence its output distribution are identical in
the two worlds.  At every $x\in(\theta_0,\theta_1]$ the true labels are
opposite.  For any randomized prediction, its two error probabilities
there sum to one.  Integrate over the interval and use nonnegativity of
the remaining errors.  The maximum is at least half the sum.
\end{proof}

\paragraph{Consequences and prior art [PROVED; KNOWN/PRIOR ART].}
OLD-only evidence cannot universally select the realized new shift or
the oracle affine weight $\Delta^2/(\Delta^2+v)$.  Identical OLD evidence
can accompany $\Delta=0$ or a very large shift.  A useful selection rule
therefore needs a declared relationship between campaigns, not merely
more OLD resampling.  The example changes $P(Y\mid X)$ and is
conditional/concept shift, not pure covariate shift.  Broader results are
given by Ben-David, Lu, Luu and Pal, AISTATS 2010,
\href{https://proceedings.mlr.press/v9/david10a.html}{Impossibility Theorems for Domain Adaptation}.

\paragraph{Experiment row [CONJECTURE: test proposal].}
Construct paired target-threshold worlds sharing exactly the same OLD
data and target score pool.  Verify identical model-selection inputs
and record the slack $R_0+R_1-(\theta_1-\theta_0)$.  A negative slack
falsifies the implementation.  Distinguishing the worlds through
additional information tests a different assumption set.

\begin{proposition}[PROVED: illustrative calibration with genuine historical campaigns]
Suppose an external anchor $\theta_0$ is known and $G$ independent
historical campaign estimates satisfy
\[
 Z_g-\theta_0\sim N(0,\tau^2+v_0),
\]
with known common $v_0$.  If $S=\sum_g(Z_g-\theta_0)^2$ and
$q_{\alpha,G}$ is the lower $\alpha$ quantile of a chi-square law with
$G$ degrees of freedom, then
\[
 \tau_U^2=\max\{0,S/q_{\alpha,G}-v_0\}
\]
is a one-sided confidence upper bound for $\tau^2$ with coverage at
least $1-\alpha$.
\end{proposition}
\begin{proof}
$S/(\tau^2+v_0)$ is chi-square with $G$ degrees of freedom.  With
probability $1-\alpha$, this ratio exceeds $q_{\alpha,G}$, implying
$\tau^2\leq S/q_{\alpha,G}-v_0$.  Truncating a negative upper expression
at zero cannot reduce coverage for the nonnegative parameter.
\end{proof}

\paragraph{Calibration limits [PROVED; protocol options CONJECTURE].}
Estimating the anchor from the same campaigns requires modified degrees
of freedom and future-prediction uncertainty; the displayed formula
does not then apply unchanged.  Random row splits from one OLD campaign
cannot identify between-campaign shift variance without further
assumptions.  Genuine historical campaigns plus exchangeability could
support nested leave-one-campaign-out calibration.  With one campaign,
a physically justified frozen shift range or transparent prespecified
sensitivity grid is defensible, but choosing its radius from NEW outcomes
is development.  A frozen rule may adapt to newly paid target labels
with their cost charged; choosing the rule after inspecting its target
endpoint is different.  NEW-136 is already open development data, so a
new transfer claim needs a separate future campaign.

\paragraph{Related robustness literature [KNOWN/PRIOR ART].}
Go and Isaac (UAI 2022),
\href{https://proceedings.mlr.press/v180/go22a.html}{Robust expected information gain using ambiguity sets},
construct an explicitly robust design objective using a KL ambiguity set
and an affine relaxation.  Sloman, Oppenheimer, Broomell and Shalizi,
\href{https://arxiv.org/abs/2205.13698}{arXiv:2205.13698},
study model misspecification and active-learning bias.  Barlas, Sloman
and Kaski,
\href{https://arxiv.org/abs/2511.07671}{arXiv:2511.07671},
develop generalized-Bayes experimental design for specified
misspecification settings.  None establishes that arbitrary logistic
coefficient regularization implements robust Bayes, or that information
gain improvements imply better Hamming acquisition in this repository.


% Round 3 geometry/discovery fragment. No preamble; all results proved below.
% Environments theorem, proposition, lemma, counterexample, proof assumed.
% Labels use the prefix rthreegeo to avoid collisions.

\section{Geometric Bayes risk and explicit regularity bridges}

\begin{proposition}[PROVED: Coherent geometric value of information]
\label{rthreegeo:coherence}
Let $G$ be a random geometric target, $\mathcal A$ a fixed action space,
and $L(a,G)\ge0$ an integrable loss. Given data $D$, define
\[
 R(D)=\inf_{a\in\mathcal A}\E[L(a,G)\mid D],\qquad
 \Delta_j(D)=R(D)-\E[R(D,O_j)\mid D].
\]
Under coherent conditioning, $\Delta_j(D)\ge0$ whenever the conditional
Bayes risks and decisions are measurable. This is an expected, one-step
statement, not an outcome-wise or multistep-optimality statement.
\end{proposition}
\begin{proof}
Retain a current Bayes action $a_D$ after every possible observation.
Optimality of the updated action gives
$R(D,O_j)\le\E[L(a_D,G)\mid D,O_j]$.
Take the conditional expectation and use the tower identity. If the
infimum is not attained, apply this argument to measurable
$\epsilon$-optimal actions and let $\epsilon\downarrow0$.
\end{proof}

\begin{proposition}[PROVED: Graph volume, medians, means and binary observations]
\label{rthreegeo:graphloss}
Let $U\subset\mathbb R^m$, let $g,h:U\to[a,b]$ be measurable, and put
$A_g=\{(u,z):u\in U,\ a\le z\le g(u)\}$.
For a nonnegative density $q$,
\[
 \mu_q(A_g\triangle A_h)
 =\int_U\left|\int_{g(u)}^{h(u)}q(u,z)\,dz\right|du.
\]
If $q(u,z)=w(u)$, this is $\int_Uw(u)|g(u)-h(u)|du$.
If $0<c\le q\le C$ between the graphs, it lies between
$c\|g-h\|_1$ and $C\|g-h\|_1$.

For a random graph $g$, the Bayes action under integrated absolute
displacement is a measurable pointwise posterior median. Under
integrated squared displacement, it is the posterior mean, and the
exact value of observing $O_j$ is
\[
 \Delta_{2,j}
 =\int_Uw(u)\operatorname{Var}\bigl(\E[g(u)\mid D,O_j]\mid D\bigr)\,du.
\]
For a binary observation $O_j\in\{-1,1\}$ with positive variance,
\[
 \Delta_{2,j}
 =\int_Uw(u)
 \frac{\operatorname{Cov}(g(u),O_j\mid D)^2}
      {\operatorname{Var}(O_j\mid D)}\,du.
\]
If every posterior graph is $L$-Lipschitz, the pointwise posterior mean
and a consistently chosen quantile median are also $L$-Lipschitz.
\end{proposition}
\begin{proof}
At each $u$, the symmetric difference of the vertical sections is exactly
the interval between $g(u)$ and $h(u)$. Tonelli's theorem proves the
identity; the density bounds follow pointwise.
Tonelli also separates each integrated posterior loss by $u$.
For a scalar $Z$, a number $a$ minimizes $\E|a-Z|$ precisely when
$\mathbb P(Z\le a)\ge1/2$ and $\mathbb P(Z\ge a)\ge1/2$; this follows
from the one-sided derivatives of the convex objective. The squared-loss
claim follows by completing the square.
Conditional total variance proves the displayed value identity.
Every square-integrable function of a nondegenerate binary variable is
affine in that variable; its conditional-mean regression coefficient is
$\operatorname{Cov}(g(u),O_j)/\operatorname{Var}(O_j)$, giving the final
formula. A deterministic observation has value zero.
Finally, the almost-sure inequality
$|g(u)-g(v)|\le L\|u-v\|$ implies the same inequality for corresponding
quantiles of the two marginal laws. Taking expectations gives it for
the means. These actions are measurable under the usual measurable
posterior-process assumptions.
\end{proof}

\begin{proposition}[PROVED: A restricted Gaussian graph problem]
\label{rthreegeo:gaussiangraph}
Suppose the current posterior of $g$ is Gaussian, with variance $s(u)^2$
and covariance $c(u,v)$, and a paid observation directly reveals $g(v)$.
Assume the displayed integrals are finite. For integrated absolute
displacement of unrestricted graph heights,
\[
 R_1=\sqrt{\frac2\pi}\int_U w(u)s(u)\,du,
\]
\[
 \Delta_1(v)=\sqrt{\frac2\pi}\int_Uw(u)
 \left[s(u)-\sqrt{s(u)^2-\frac{c(u,v)^2}{s(v)^2}}\right]du.
\]
For integrated squared displacement,
\[
 \Delta_2(v)=\int_Uw(u)\frac{c(u,v)^2}{s(v)^2}\,du.
\]
The values are zero if $s(v)=0$. A binary membership observation does not
directly reveal $g(v)$ and is a different observation experiment.
\end{proposition}
\begin{proof}
A Gaussian marginal's mean is a median and its mean absolute centered
deviation is $\sqrt{2/\pi}s(u)$. Gaussian conditioning changes the variance
to $s(u)^2-c(u,v)^2/s(v)^2$, independently of the observed value.
Substitute this into the absolute-deviation risk. For squared loss,
subtract the integrated updated variance from the integrated current
variance. A zero-variance observation is already known.
\end{proof}

\begin{proposition}[PROVED: Graph Hausdorff and density-dependent $L^1$ control]
\label{rthreegeo:holderbridge}
Let $g,h$ be $L$-Lipschitz on the same compact $U$, and let
$\Gamma_g=\{(u,g(u)):u\in U\}$. Then
\[
 \frac{\|g-h\|_\infty}{\sqrt{1+L^2}}
 \le H(\Gamma_g,\Gamma_h)\le\|g-h\|_\infty.
\]
More specifically, let $U=[0,1]^m$, let $f=g-h$ be $K$-Lipschitz,
$K>0$, and let $w\ge w_{\min}>0$. Put $M=\|f\|_\infty$ and
$r=\min\{1/2,M/(2K\sqrt m)\}$. Then
\[
 \int_Uw|f|\ge w_{\min}\frac M2 r^m.
\]
In the regime $M\le K\sqrt m$ this gives
\[
 M\le\left[
 \frac{2^{m+1}K^m m^{m/2}}{w_{\min}}\int_Uw|f|
 \right]^{1/(m+1)}.
\]
\end{proposition}
\begin{proof}
Match graph points with the same base coordinate for the upper bound.
At $u$ maximizing $|g(u)-h(u)|=M$, a point $(v,h(v))$ is at least
\[
 \sqrt{t^2+(M-Lt)_+^2},\qquad t=\|u-v\|,
\]
away. The minimum over $t\ge0$ is $M/\sqrt{1+L^2}$.
For the integral bound, at a maximizer of $|f|$ choose, in each coordinate,
a direction with at least $1/2$ available space inside the unit cube.
The resulting one-sided box of side $r$ is contained in $U$. Every point
in it lies within $\sqrt m r$ of the maximizer, and therefore has
$|f|\ge M/2$. Integrate on this box and rearrange. For $K=0$, $f$ is
constant and the integral is exactly $M\int w$. A cone of height $M$
and slope $K$ has mass of order $M^{m+1}/K^m$, showing that the
small-error exponent cannot generally be improved under only a
Lipschitz bound.
\end{proof}

\begin{proposition}[PROVED: A finite graph approximation]
\label{rthreegeo:quadrature}
Let a probability measure $\nu$ on $U$ be partitioned into cells $C_i$
of diameter at most $h$, with $u_i\in C_i$ and $w_i=\nu(C_i)$.
For graphs with Lipschitz constants $L_g,L_h$,
\[
 \left|\int_U|g-h|\,d\nu-
 \sum_iw_i|g(u_i)-h(u_i)|\right|\le(L_g+L_h)h.
\]
If the base points form an $h$-net, the difference between the Hausdorff
distance of the sampled graph clouds and that of the complete graphs
is at most
\[
 h\left(\sqrt{1+L_g^2}+\sqrt{1+L_h^2}\right).
\]
\end{proposition}
\begin{proof}
The function $|g-h|$ is $(L_g+L_h)$-Lipschitz. Integrate its deviation
from each cell representative. Every graph point is within
$h\sqrt{1+L_g^2}$ of its sampled graph cloud, and conversely every cloud
point lies on the graph. Apply the triangle inequality for Hausdorff
distance twice. These arguments require graph heights; a finite label
pool alone need not locate an intervening root.
\end{proof}

\begin{proposition}[PROVED: A density-weighted normal-displacement expansion]
\label{rthreegeo:tube}
Let $S$ be a compact oriented $C^2$ hypersurface of dimension $m$ with
injective normal coordinates $F(s,t)=s+tn(s)$ for $|t|<r$.
Suppose its principal curvatures are bounded in magnitude by $K$.
Let $0<a<r$, $aK<1$, and assume two regions differ exactly by the
normal strip between $t=0$ and $t=\delta(s)$, where $|\delta|\le a$;
there are no additional components or orientation changes.
Let $q\ge0$ be bounded by $Q$ and $L_q$-Lipschitz on the tube. Then
\[
 \mu_q(A\triangle B)=
 \int_S\left|\int_0^{\delta(s)}q(F(s,t))J(s,t)\,dt\right|dA(s),
 \quad J(s,t)=\prod_{i=1}^m(1-t\kappa_i(s)),
\]
and
\[
 \left|\mu_q(A\triangle B)-\int_Sq(s)|\delta(s)|\,dA(s)\right|
 \le\frac C2\int_S\delta(s)^2\,dA(s),
\]
where
\[
 C=L_q(1+aK)^m+QmK(1+aK)^{m-1}.
\]
\end{proposition}
\begin{proof}
The normal-coordinate change of variables gives the first identity.
The assumptions make its Jacobian positive. Telescoping the product
gives
$|J-1|\le mK|t|(1+aK)^{m-1}$ and $|J|\le(1+aK)^m$.
Consequently
$|q(F(s,t))J(s,t)-q(s)|\le C|t|$.
Integrating from $0$ to $\delta(s)$ and then over $S$ proves the bound.
\end{proof}

\begin{counterexample}[COUNTEREXAMPLE; PROVED: Latent amplitude does not identify boundary uncertainty]
\label{rthreegeo:scale}
Let $\Theta\sim N(m,\sigma^2)$ and $f_c(x)=c(x-\Theta)$, $c>0$.
The boundary is $\{\Theta\}$ for all $c$, although
$\operatorname{sd}(f_c(x))=c\sigma\to0$ as $c\downarrow0$.
Squared boundary Bayes risk remains $\sigma^2$; the scale-invariant
quantity $\operatorname{sd}(f_c)/|f_c'|$ equals $\sigma$.
A noiseless sign query at $x=m$ has squared-boundary value
$2\sigma^2/\pi$. Under the different observation contract
$O=\operatorname{sign}(f_c(m)+\eta)$ with independent
$\eta\sim N(0,\tau^2)$, the value is
\[
 \frac{2c^2\sigma^4}{\pi(c^2\sigma^2+\tau^2)}.
\]
\end{counterexample}
\begin{proof}
The root and scale claims follow directly from the definition.
A noiseless query splits the Gaussian threshold at its median; the two
conditional means are $m\pm\sigma\sqrt{2/\pi}$, whose variance is
$2\sigma^2/\pi$. More generally put $X=\Theta-m$ and
$Z=-cX+\eta$. Joint Gaussian regression gives
\[
 \E[X\mid\operatorname{sign}Z]
 =\frac{\operatorname{Cov}(X,Z)}{\operatorname{sd}(Z)}
   \sqrt{\frac2\pi}\operatorname{sign}Z.
\]
Its variance is the displayed value by conditional total variance.
The fixed-noise value tends to zero while the random boundary law stays
unchanged.
\end{proof}

\begin{proposition}[PROVED: Local root sensitivity requires a nonzero gradient]
\label{rthreegeo:rootsensitivity}
On a fixed normal line, let $f_t(x)=f(x)+t e(x)$ be continuously
differentiable near $(t,x)=(0,x_0)$, with $f(x_0)=0$ and
$f'(x_0)\ne0$. There is a local differentiable root $x(t)$, and
\[
 x'(0)=-\frac{e(x_0)}{f'(x_0)}.
\]
This derivative is invariant when both $f$ and $e$ are multiplied by a
positive constant.
\end{proposition}
\begin{proof}
The implicit function theorem gives the root and its differentiability.
Differentiate $f(x(t))+t e(x(t))=0$ at zero and solve for $x'(0)$.
The multiplicative constant cancels. Uniform positional bounds require
a uniform derivative lower bound and control of additional roots;
the local identity by itself supplies neither.
\end{proof}

\begin{counterexample}[COUNTEREXAMPLE; PROVED: Area and tolerance overlap are different losses]
\label{rthreegeo:areansd}
For $h=0$ and $g_k(u)=k^{-1}\sin(k^2u)$ on $[0,1]$, graph Hausdorff
distance and symmetric-difference area tend to zero, whereas graph
length diverges. For two parallel equal-area surfaces separated by $d$,
normalized surface Dice at tolerance $\tau$ is one for $d\le\tau$ and
zero for $d>\tau$, while Hausdorff distance is $d$.
\end{counterexample}
\begin{proof}
Uniform graph displacement is at most $1/k$, bounding both Hausdorff
distance and integrated displacement. Graph length is at least
$\int_0^1 k|\cos(k^2u)|du=(2/\pi)k+o(k)$.
For the parallel surfaces every point's nearest-surface distance is
exactly $d$, so every point either counts within tolerance or does not.
\end{proof}

\section{A Hausdorff acquisition counterexample beyond pairwise moments}

\begin{lemma}[PROVED: Bayes sup-distance risk on a binary cube]
\label{rthreegeo:cuberisk}
For any distribution of $B\in\{0,1\}^m$, allowing all estimates
$h\in[0,1]^m$,
\[
 \inf_h\E\|h-B\|_\infty
 =\min\left\{\frac12,1-\max_b\mathbb P(B=b)\right\}.
\]
\end{lemma}
\begin{proof}
The cube center has risk $1/2$, and a modal vertex has risk $1-p_{\max}$.
For any $h$, round each coordinate to a nearest endpoint to obtain $v$,
and let $r=\|h-v\|_\infty\le1/2$.
Every other vertex is at distance at least $1-r$, hence the risk is at
least $p_vr+(1-p_v)(1-r)$. Its minimum over $r\in[0,1/2]$ is
$\min\{1/2,1-p_v\}$, which is at least the stated lower bound.
The two explicit actions attain that bound.
\end{proof}

\begin{counterexample}[COUNTEREXAMPLE; PROVED: Equal binary pair moments, different geometric query values]
\label{rthreegeo:parity}
Let three bases $u_i$ have pairwise distances at least $3$.
A random boundary is $\{(u_i,B_i):i=1,2,3\}$, where
$B\in\{0,1\}^3$, and admissible estimates have heights in $[0,1]^3$.
Querying column $i$ at height $1/2$ reveals $B_i$.
The following two laws have identical first and pairwise binary moments:
\[
\begin{array}{c|rrrrrrrr}
 B&000&001&010&011&100&101&110&111\\\hline
 70P_+(B)&34&0&0&12&0&12&12&0\\
 70P_-(B)&22&12&12&0&12&0&0&12
\end{array}
\]
Their current Hausdorff Bayes risks are both $1/2$, but the exact values
of querying $B_1$ are $11/70$ and zero, respectively.
\end{counterexample}
\begin{proof}
Because other columns are at distance at least $3$ whereas the
same-column height difference is at most $1$, the Hausdorff loss is
exactly $\|h-B\|_\infty$.
Direct summation gives, for both laws,
\[
 \mathbb P(B_i=1)=24/70,\qquad
 \mathbb P(B_i=B_j=1)=12/70\quad(i\ne j).
\]
These determine every first and pairwise binary moment.
Every atom has probability less than $1/2$, so
Lemma~\ref{rthreegeo:cuberisk} gives current risk $35/70$.
The events $B_1=0$ and $B_1=1$ have masses $46/70$ and $24/70$.
For a branch of mass $s$ and largest atom mass $p$, its unnormalized
Bayes risk is $\min\{s/2,s-p\}$ by the same lemma.
Under $P_+$, branch zero has largest atom $34/70$ and risk $12/70$;
branch one has risk $12/70$. Updated expected risk is $24/70$.
Under $P_-$, branch zero has largest atom $22/70$, so its risk is
$23/70$; branch one again has risk $12/70$. Updated expected risk
is $35/70$. Subtraction proves the two query values.
This counterexample permits intermediate height estimates and thus
does not rely on forcing the estimate to be a posterior sample.
\end{proof}

\section{Discovery, refinement, and violations of structural assumptions}

\begin{counterexample}[COUNTEREXAMPLE; PROVED: Optimal startup, suboptimal later geometry at equal budget]
\label{rthreegeo:discoverygeometry}
Fix any $k\ge1$, put $M=k+1$, and let
$\Theta_1,\ldots,\Theta_M$ be independent fair bits.
The geometric target has $M$ separated columns, heights $H\Theta_i$,
and base separations at least $3H$.
The finite candidate pool contains one membership query $q_i$ at height
$H/2$ in each column, revealing $\Theta_i$, and two paid certification
anchors of fixed labels $0$ and $1$.
Startup requires actually observed labels of both classes.

Policy A queries the anchors first. It has the globally minimal
$T_{\rm both}=2$ surely. With an optimal continuation and total paid
budget $Q=k+2$, its final Hausdorff Bayes risk is $H/2$.
Policy B has $\E T_{\rm both}=5/2$, $T_{\rm both}\le3$, and zero final
geometric risk by the same total budget $Q$.
\end{counterexample}
\begin{proof}
No policy can observe both labels in fewer than two queries, proving
A's startup optimality. After its two anchors, A has at most $k$
individual-bit queries for $M=k+1$ independent bits.
Conditional on any adaptive query history, an unqueried bit remains
independent and fair. Its minimum expected absolute height error is
$H/2$, a lower bound on expected Hausdorff error.
Set every unknown height to $H/2$ and every known height correctly;
the loss is exactly $H/2$, attaining the bound.

B first queries $q_1,q_2$. With probability $1/2$ their labels differ,
so startup ends at two. Otherwise it queries the opposite-label anchor,
ending at three. Thus its expected startup cost is $5/2$.
It then queries the remaining $M-2=k-1$ bits. The total cost is at most
$3+(k-1)=Q$, and all boundary heights are known. Unused budget can be
padded by an irrelevant remaining query if an exactly equal paid count
is desired. Therefore its risk is zero.

This is an operational paid-certification example. If known structural
anchors may replace actual paid observations, the startup contract
changes. The result concerns the stated actual-observation contract.
\end{proof}

\begin{proposition}[PROVED: State coalescence bounds later effects]
\label{rthreegeo:coupling}
Couple two policies on the same truth and continuation seed. Suppose
their complete sufficient states coincide at budget $b$ with probability
at least $1-\delta$, and they use the same subsequent transition and
prediction rules. For any later endpoint in an interval of length $L$,
\[
 |\E Z_A-\E Z_B|\le L\delta.
\]
For a weighted learning-curve endpoint with nonnegative weights
$\lambda_t$ summing to one, if $\delta_t$ bounds the probability of
noncoalescence at budget $t$, then
\[
 |\E\textstyle\sum_t\lambda_t Z_{A,t}
       -\E\sum_t\lambda_t Z_{B,t}|
 \le L\sum_t\lambda_t\delta_t.
\]
Identical query sets imply identical states only under order-invariant
fitting, equal seeds, and absence of additional history-dependent state.
\end{proposition}
\begin{proof}
Identical transitions preserve identical states inductively on the
coalescence event. Endpoint differences vanish on that event and are
bounded by $L$ on its complement. Take expectations.
Apply this argument at each budget and use the triangle inequality for
the weighted endpoint.
\end{proof}

\begin{counterexample}[COUNTEREXAMPLE; PROVED: Three different meanings of monotonicity]
\label{rthreegeo:monotone}
Monotone class probabilities do not permit deterministic label
propagation: independent labels with probabilities $0.49,0.51$ are
stochastically ordered but have probability $0.49^2$ of the reversed
pair $(1,0)$.

On a chain of $N$ points, the deterministic truth
$(1,0,\ldots,0,1)$ differs at one point from the monotone threshold
$(0,\ldots,0,1)$. Propagating the first observed positive upward gives
$N-2$ wrong labels. This is one contaminated point but $N-2$ violating
pairs. Even exactly one violating pair destroys zero-error propagation:
the sequence $(0,\ldots,0,1,0,1,\ldots,1)$ with adjacent $1,0$
incorrectly propagates the first of those two labels to the second.
\end{counterexample}
\begin{proof}
Independence gives $\mathbb P(Y_1=1,Y_2=0)=0.49(1-0.51)$.
For the first deterministic sequence, every middle zero is dominated
by the first positive and is therefore incorrectly removed by the
hard monotonicity rule. Each is also a distinct violating pair with
the first point. In the second sequence the adjacent reversed pair is
the only inversion; propagation across it is nonetheless wrong.
\end{proof}

\begin{proposition}[PROVED: A robust and tight rank-error discovery bound]
\label{rthreegeo:rank}
In a fixed tie-broken order of $N$ items containing $r\ge1$ designated
rare items, let $V$ count common-before-rare pairs, and let $D$ be the
first rare rank. Then
\[
 D\le1+\left\lfloor\frac Vr\right\rfloor.
\]
If $0<r<N$ and AUC is computed on this strict rare-oriented ranking,
\[
 D\le1+\lfloor(1-\mathrm{AUC})(N-r)\rfloor.
\]
For a top-$M$ subset containing $s$ rare items,
\[
 V\ge(M-s)(r-s).
\]
In particular $V<Mr$ guarantees that the first $M$ positions contain
a rare item. The first-hit bound is attained when $V=kr$ by placing
$k$ common items, all $r$ rare items, then all remaining common items.
\end{proposition}
\begin{proof}
Each of the $D-1$ common items preceding the first rare precedes all
$r$ rare items. Thus $V\ge r(D-1)$.
For a strict ranking, $1-\mathrm{AUC}=V/[r(N-r)]$.
For the tail inequality, each of the $M-s$ common items inside the
top-$M$ subset precedes every one of the $r-s$ rare items outside it.
If $s=0$, this forces $V\ge Mr$.
The displayed block ordering has exactly $kr$ inversions and first-hit
rank $k+1$, proving attainment. Half-credit AUC for unresolved ties
does not fix a deterministic first-hit rank; the declared tie order
must be used.
\end{proof}


% Root synthesis results, complete proofs; no external file dependencies.
\section{Additional decision certificates and finite-pool checks}

\begin{proposition}[PROVED: exact zero-value and limiting conditions]
\label{prop:zero-value-r3}
If Bayes actions are attained under a fixed loss, coherent Bayes value
is zero exactly when some current Bayes action is conditionally optimal
after almost every observation. For a sequence of binary experiments,
$P(O_j=1)\to1/2$ is equivalent to $\mathbb E O_j\to0$.
For fixed finite Hamming targets with positive weights, query value tends
to zero if and only if
$(|\mathbb E[T_iO_j]|-|\mathbb E T_i|)_+\to0$ for every target.
For integrated squared graph displacement with a binary observation
whose mean tends to zero, value tends to zero if and only if
$\int w(u)\operatorname{Cov}(g(u),O_j)^2du\to0$.
\end{proposition}
\begin{proof}
For a current Bayes action $a_0$, the value equals the expectation of
the nonnegative conditional optimality gap
$\mathbb E[L(a_0,T)\mid O]-\inf_a\mathbb E[L(a,T)\mid O]$.
It vanishes exactly when that gap vanishes almost surely.
The observation-marginal equivalence follows from
$P(O=1)=(1+\mathbb E O)/2$. The Hamming result follows from the
binary-observation identity and a finite sum of nonnegative terms
with fixed positive weights. The squared-graph formula is
Proposition~\ref{rthreegeo:graphloss}; its denominator
$\operatorname{Var}(O)=1-(\mathbb E O)^2$ tends to one.
\end{proof}

\begin{proposition}[PROVED: a self-information condition for observed-label margin]
\label{prop:self-margin-r3}
Let fixed binary targets have nonnegative weights summing to one. Put
$b_j=\min\{P(O_j=1),P(O_j=-1)\}$. Then their total Hamming Bayes
value $V(j)$ is at most $b_j$. Suppose every candidate $j$ is also a target
with weight $w_j>0$, and its unweighted self-value $v_{jj}$ is at least
$\alpha b_j$ for a common $\alpha>0$. If $m$ maximizes $b_j$, then
\[
 V(m)\ge \alpha w_m\max_j V(j)\ge\alpha w_{\min}\max_j V(j).
\]
It suffices that the self-value condition hold at the selected $m$.
\end{proposition}
\begin{proof}
Let $o_+$ be a most probable observation outcome and $d_+,d_-$ the
conditional Bayes actions. Before querying, using $d_+$ is feasible,
so the old Bayes risk is at most its unconditional risk. Subtracting the
optimal observation-dependent risk leaves at most
\[
 P(O_j=o_-)\,
 \mathbb E[L(d_+,T)-L(d_-,T)\mid O_j=o_-]\le b_j,
\]
because normalized Hamming loss lies in $[0,1]$.
The selected query's total value includes its nonnegative self contribution,
so $V(m)\ge w_m v_{mm}\ge\alpha w_m b_m$. Finally
$b_m\ge b_j\ge V(j)$ for each $j$.
\end{proof}

\begin{counterexample}[COUNTEREXAMPLE: noiseless disjoint targets give no margin ratio]
Let the sole target $T$ have $P(T=1)=1/2+\delta$, where $0<\delta<1/2$.
Query $a$ reveals an independent fair sign; query $b$ reveals $T$.
Observed-label margin uniquely selects $a$, whose value is zero.
Query $b$ has value $1/2-\delta$.
\end{counterexample}
\begin{proof}
The fair observation is strictly more uncertain than the biased one.
Independence leaves the target posterior unchanged after $a$; revelation
after $b$ removes all target error.
\end{proof}

\begin{proposition}[PROVED: model-internal improvement and updater error]
\label{prop:improvement-r3}
In Theorem~\ref{thm:terminal-transfer-r3}'s notation, let
$\gamma=V_Q(j_Q)-V_Q(j_m)$, and use the $Q$-Bayes terminal rules
$g_{j_Q},g_{j_m}$. The true-risk improvement from replacing $j_m$ by
$j_Q$ is at least
\[
 \gamma-\frac{D_{j_Q}+D_{j_m}}2.
\]
If actual terminal rules $\widetilde g_j$ differ from $g_j$, define
\[
 \zeta_j=\sum_iw_iP\{\widetilde g_{ij}(O_j)\ne g_{ij}(O_j)\}.
\]
The same improvement lower bound for the actual rules subtracts the
additional term $\zeta_{j_Q}+\zeta_{j_m}$.
\end{proposition}
\begin{proof}
The fixed-rule risk discrepancy is at most $D_j/2$ by
Theorem~\ref{thm:terminal-transfer-r3}. Insert the two corresponding
$Q$ risks between the true risks. Their difference is $\gamma$ because
the baseline cancels. A classification loss can change only when the
two predicted labels differ; the absolute risk difference caused by
replacing $g_j$ is at most $\zeta_j$. Apply the triangle inequality.
\end{proof}

\begin{proposition}[PROVED: covariance and joint-log-score versions]
\label{prop:cov-kl-r3}
Write $n_j^R=\mathbb E_R O_j$ and
$k_{ij}^R=\operatorname{Cov}_R(T_i,O_j)$. If both target and observation
spin marginals differ by at most $\epsilon_m$, and
$\sum_iw_i|k_{ij}^P-k_{ij}^Q|\le\epsilon_c$ for every $j$, then
the terminal regret in Theorem~\ref{thm:terminal-transfer-r3} is at most
$2\epsilon_m+\epsilon_c$.
Alternatively, if $P_{ij},Q_{ij}$ are the actual joint laws of
$(T_i,O_j)$, then
\[
 D_j\le\sum_iw_i\sqrt{2\operatorname{KL}(P_{ij}\Vert Q_{ij})}
 \le\sqrt{2\sum_iw_i\operatorname{KL}(P_{ij}\Vert Q_{ij})}.
\]
Thus a uniform bound $K$ on the last weighted KL sum gives terminal
regret at most $\sqrt{2K}$.
\end{proposition}
\begin{proof}
Since $c_{ij}=m_i n_j+k_{ij}$ and spin means are bounded by one,
$|c_{ij}^P-c_{ij}^Q|\le |m_i^P-m_i^Q|+|n_j^P-n_j^Q|
+|k_{ij}^P-k_{ij}^Q|$. The same bound dominates the target-mean
error, proving $D_j\le2\epsilon_m+\epsilon_c$.
Each of the two moments defining $D_j$ is an expectation of a
function in $[-1,1]$, so its difference is at most $2\operatorname{TV}
(P_{ij},Q_{ij})\le\sqrt{2\operatorname{KL}(P_{ij}\Vert Q_{ij})}$
by Pinsker. Jensen's inequality for the square root gives the weighted
bound. Expected excess negative log score equals KL by its definition;
raw negative log score also contains the true entropy and is not KL.
\end{proof}

\begin{proposition}[PROVED: vanishing information under a fixed smooth channel]
\label{prop:bounded-scale-r3}
Let $G$ be a random state/field and let a query return a binary observation
with $P_c(O=1\mid G)=\kappa(cG_j)$, where $\kappa(0)=1/2$.
For any fixed measurable loss in $[0,1]$, its coherent Bayes information
value satisfies
\[
 0\le V_c(j)\le\mathbb E|\kappa(cG_j)-1/2|.
\]
Consequently the value vanishes as $c\downarrow0$ for a continuous
bounded channel. If $\mathbb E|G_j|<\infty$, logistic gives the upper
bound $c\mathbb E|G_j|/4$, and $\kappa(f)=\Phi(f/\tau)$ gives
$c\mathbb E|G_j|/(\tau\sqrt{2\pi})$.
\end{proposition}
\begin{proof}
Compare the joint law of $(G,O)$ with the law sharing the same $G$
but making $O$ an independent fair coin. Their total variation is
$\mathbb E|\kappa(cG_j)-1/2|$, as follows by summing the two outcome
probability differences conditionally on $G$. The risk of any fixed
observation-to-action rule changes by at most this TV, and taking the
infimum preserves the same bound. Under the independent coin no rule
improves on the prequery Bayes action, so its optimal risk is the
prequery risk. Nonnegativity follows from retaining that action.
Dominated convergence and the respective derivative bounds on the
links finish the proof.
\end{proof}

\begin{proposition}[PROVED: independent copy blocks make EER optimal for every budget]
\label{prop:blocks-r3}
Partition a finite target pool into blocks. Labels are identical within
each block and independent between blocks. A noiseless query reveals
the block's common sign. Let $W_b$ be total block weight and $u_b$
its initial Bayes error. Optimal total gain at budget $B$ is the sum
of the largest $\min(B,\text{number of blocks})$ numbers $W_bu_b$.
Exact greedy EER achieves it.
\end{proposition}
\begin{proof}
A query resolves its block and leaves every other block law unchanged.
Hence the value of the first query in block $b$ is always $W_bu_b$,
and additional queries there have zero value. Every adaptive
realization can resolve at most $B$ blocks, so its gain is bounded
by the sum of the largest fixed block values. Choosing those blocks
attains the bound, independently of the observed signs.
\end{proof}

\begin{counterexample}[COUNTEREXAMPLE: correct-specification multistep EER exceeds margin]
Take two independent fair singleton blocks and an independent
$M$-point copy block of error $1/2-\eta$, with equal point weights
$1/(M+2)$ and $0<\eta<1/2$. If $M(1/2-\eta)>1/2$, then at budget two
margin gains $1/(M+2)$, whereas exact EER gains
$[M(1/2-\eta)+1/2]/(M+2)$ and is horizon-optimal.
\end{counterexample}
\begin{proof}
The two fair singleton marginals are strictly most uncertain until
both are queried, so margin spends both queries there. Each contributes
$1/[2(M+2)]$. The copy-block value is strictly larger than either
singleton value; the preceding proposition gives the EER ordering
and optimum.
\end{proof}

\begin{proposition}[PROVED: randomization cannot improve a known finite-policy optimum]
For a fixed known finite-world law, finite actions/outcomes, a fixed
horizon and terminal loss, an optimal adaptive policy can be chosen
deterministic.
\end{proposition}
\begin{proof}
Condition on all internal random choices of a randomized decision tree.
Its expected risk is a mixture of risks of deterministic trees, so at
least one tree has risk no greater than the mixture. Equivalently,
backward induction minimizes an affine function of each action
distribution at an extreme point. This does not address minimax
robustness over an unknown law, where mixtures can help.
\end{proof}

\begin{proposition}[PROVED: edge-loss Bayes actions and the third moment]
\label{prop:edge-moment-r3}
On a fixed weighted graph, let $Z_e=T_iT_k$ denote whether endpoint
signs agree. For unconstrained edge-sign predictions $a_e\in\{-1,1\}$,
weighted edge-disagreement Bayes risk is
\[
 R_E=\frac12\sum_e w_e(1-|\mathbb E Z_e|).
\]
The exact one-step reduction after binary $O_j$ is
\[
 \frac12\sum_e w_e
 (|\mathbb E[T_iT_kO_j]|-|\mathbb E[T_iT_k]|)_+.
\]
If the output must be a cut induced by vertex predictions, the
optimization is over the feasible cut signs instead. Pair and triple
moments remain sufficient for evaluating and optimizing this fixed
loss before and after one binary observation, but coordinatewise edge
modes need not form a feasible cut.
\end{proposition}
\begin{proof}
Apply Lemma~\ref{lem:binary-observation-r3} to the binary target $Z_e$.
For constrained actions, the expected loss of any vertex labeling is
an affine combination of the edge means, so minimizing over finitely
many labelings uses only those means. Conditional edge means are
\[
 \mathbb E[Z_e\mid O_j=o]
 =\frac{\mathbb E Z_e+o\mathbb E[Z_eO_j]}
        {1+o\mathbb E O_j}
\]
on positive-probability outcomes. Thus second and third moments plus
the observation marginal suffice. Cycles impose compatibility
constraints: an odd number of cut edges on a cycle is impossible,
illustrating why independent edge modes may fail feasibility.
\end{proof}

\begin{counterexample}[COUNTEREXAMPLE: even an optimized ratio can have negative coherent value]
\label{ce:optimized-ratio-r3}
Use two disjoint graph edges with cut indicators $(1,Z)$, where
$P(Z=1)=p$. All four cut predictions are feasible. Consider the
posterior functional
\[
 R(p)=\min_{a\in\{0,1\}^2}
 \frac{\mathbb E\sum_{e=1}^2|Z_e-a_e|}
      {\mathbb E\sum_{e=1}^2 Z_e+\sum_{e=1}^2a_e}.
\]
At $p=3/4$, $R(p)=1/15$. A coherent binary observation with equally
probable posterior probabilities $p'=1/2$ and $p'=1$ gives expected
future risk $1/14$, hence reduction $-1/210$.
\end{counterexample}
\begin{proof}
For predictions $(0,0),(1,0),(0,1),(1,1)$ the respective ratios are
\[
 1,\qquad \frac p{2+p},\qquad\frac{2-p}{2+p},
 \qquad\frac{1-p}{3+p}.
\]
For $p\in[1/2,1]$, the last is minimal. At $3/4$ it is $1/15$,
at $1/2$ it is $1/7$, and at $1$ it is zero.
Define the observation by the joint probabilities
$P(O=0,Z=0)=1/4$, $P(O=0,Z=1)=1/4$,
$P(O=1,Z=1)=1/2$. Its two branches have the stated masses and
posteriors; no incoherent update is used. The expected future ratio
is $1/14$. Unlike this ratio of expectations, the posterior expectation
of a fixed normalized loss, minimized over fixed admissible actions,
obeys Proposition~\ref{prop:nonnegative-r3}.
\end{proof}

\begin{proposition}[PROVED: uniform first-hit and both-class discovery laws]
\label{prop:tail-r3}
For $0<r<N$, a uniform random ordering of $N$ items with $r$ designated
rare items has first-hit time $D$ satisfying, for $0\le k\le N$,
\[
 P(D>k)=\frac{\binom{N-r}{k}}{\binom Nk},\qquad
 \mathbb E D=\frac{N+1}{r+1}.
\]
Its time $T$ to observe both classes satisfies, for $1\le k\le N$,
\[
 P(T>k)=\frac{\binom rk+\binom{N-r}k}{\binom Nk},\qquad
 \mathbb E T=\frac{N+1}{r+1}+\frac{N+1}{N-r+1}-1.
\]
Binomial coefficients exceeding available items are zero; survivors
are zero beyond the population size.
In a fixed subset with $M$ items and $s\ge1$ rare items, uniform
ordering has the same first-hit formulas with $N,r$ replaced by $M,s$.
If $s/M\ge r/N$, this subset first-hit time stochastically dominates
in the favorable direction: its survival probability is no larger.
\end{proposition}
\begin{proof}
No rare item among the first $k$ is precisely a size-$k$ subset chosen
from the $N-r$ common items, giving the survival formula.
The $N-r$ common items occupy $r+1$ exchangeable gaps around the rare
items, so the expected number before the first rare is $(N-r)/(r+1)$.
For $T>k$ with $k\ge1$, all $k$ labels are rare or all are common,
disjoint events. Summing their survivor probabilities gives the
displayed mean, with the duplicated $k=0$ term corrected by subtracting
one. Apply the same argument to the subset.
Conditional on $k$ failures, the subset success hazard is $s/(M-k)$
and the full-pool hazard is $r/(N-k)$ while both are defined.
The difference after cross multiplication has numerator
$sN-rM+k(r-s)\ge0$, because $s\le r$ and enrichment holds.
Multiplying failure hazards proves the survivor ordering; after
the subset's common items are exhausted its survival is zero.
Neither argument applies to an arbitrary history-conditioned
deterministic extreme-selection policy.
\end{proof}

\begin{counterexample}[COUNTEREXAMPLE: no arbitrary continuous-boundary guarantee from finite signs]
\label{ce:finite-boundary-r3}
Let a finite query pool in an interval leave an open gap $(a,b)$.
Two noiseless smooth classifiers $f_0(x)=x-\theta_0$ and
$f_1(x)=x-\theta_1$, with $a<\theta_0<\theta_1<b$, agree on every
pool label but have boundary Hausdorff distance $\theta_1-\theta_0$.
Any pool-only estimator has worst-world expected boundary error at
least $(\theta_1-\theta_0)/2$.
\end{counterexample}
\begin{proof}
Every pool point is on the same side of both thresholds. Thus even
revealing all pool labels gives identical data in the two worlds.
Couple the estimator's randomness so its output is the same.
The triangle inequality for Hausdorff distance implies that its two
errors sum to at least the distance between the boundaries. Take
expectations and the larger error. The thresholds can approach the
gap endpoints. A coverage bound limits this ambiguity but an arbitrary
finite pool supplies no smaller universal positional guarantee.
\end{proof}

\section*{Verification and use}
\addcontentsline{toc}{section}{Verification and use}
\paragraph{Verification [NUMERICALLY CHECKED].}
The companion verification record checks the binary identity on 1,819
rational laws, the exact negative edge/Dice examples, the parity
example's identical initial pairs, and 2,036 ranked binary sequences.
The terminal moment bound was additionally checked on 2,000 random
joint-law pairs. These checks supplement the proofs and do not
establish that any assumption holds for a real campaign.
\paragraph{Experiment mapping [CONJECTURE: test proposals].}
Every useful result is paired with a computable statistic and falsifier
in the report's experiment table. Synthetic algebra checks, same-state
model diagnostics, and future-campaign policy comparisons are distinct
forms of evidence. No new empirical acquisition experiment was run.
\end{document}
